\section{Introduction}
\label{sec:introduction}

The rapid shift towards containerized and cloud-native infrastructures has fundamentally transformed the landscape of runtime threat detection.
Modern security systems must now operate within highly dynamic environments, where the density of microservices and the velocity of system calls demand monitoring solutions that offer not only comprehensive visibility but also near-zero performance interference.
Despite significant progress, current state-of-the-art tools remain trapped in a tripartite trade-off between \textit{detection capability}, \textit{execution performance}, and \textit{isolation safety}.

\textbf{The Trilemma of Runtime Monitoring.}
Existing approaches to runtime security generally fall into two categories, both of which exhibit critical shortcomings.
First, rule-based systems such as Falco~\cite{falco} and Tetragon~\cite{tetragon} leverage eBPF for high-performance event filtering. While efficient, these systems are fundamentally limited to detecting known attack patterns, failing to generalize to novel or obfuscated exploits.
Second, machine learning (ML) based intrusion detection systems (IDS) offer the promise of behavioral anomaly detection. However, traditional ML pipelines—often reliant on heavy Python runtimes or complex inference engines—introduce latencies in the range of tens to hundreds of milliseconds. This is several orders of magnitude higher than the microsecond-level requirements of inline system call monitoring.
Furthermore, neither approach adequately addresses \textit{rule isolation}. In existing frameworks, detection logic typically executes with high privileges or within the main monitoring process, meaning a single compromised or malicious rule can access sensitive host resources, exfiltrate event data, or crash the entire security daemon.

\textbf{Our Insight: Synergistic Co-design.}
We argue that the perceived incompatibility between advanced ML-based detection and microsecond latency is an artifact of suboptimal implementation rather than a fundamental theoretical limit.
Our key insight is that by co-designing the entire monitoring pipeline—from kernel-level event capture to sandboxed user-space inference—we can achieve ML-grade detection without sacrificing performance or safety.
Specifically, we leverage the high-performance eBPF subsystem for event capture, a zero-copy shared memory channel for data transfer, and the WebAssembly (Wasm) runtime for secure, near-native execution of detection logic.

\textbf{WasmSentinel.}
We present \textsc{WasmSentinel}, a high-performance framework that executes arbitrary detection logic—including neural network inference—within lightweight Wasm sandboxes.
The core of \textsc{WasmSentinel} is a hybrid detection architecture that combines a sub-microsecond n-gram fast-path (STIDE) with a deeper, Wasm-optimized Multi-Layer Perceptron (MLP) for suspicious events.
Crucially, by executing all detection logic within Wasm, \textsc{WasmSentinel} provides strong hardware-level isolation and capability-based security (via WASI), enabling the safe deployment of third-party or community-contributed detection rules without trusting their authors.

\textbf{Key Results and Contributions.}
We evaluate \textsc{WasmSentinel} using the ADFA-LD dataset and modern simulated attacks, including container escapes (CVE-2024-21626) and privilege escalations (CVE-2024-1086).
Our experimental results demonstrate that \textsc{WasmSentinel} achieves:
\begin{enumerate}
    \item \textbf{High-Fidelity Detection:} An F1-score of 0.84 on ADFA-LD, outperforming complex deep learning baselines such as LSTMs and GNNs while maintaining a $45\mu$s logic latency.
    \item \textbf{Microsecond Latency:} A median inference latency of $33\mu$s, enabling real-time detection on high-throughput system call streams.
    \item \textbf{Superior Efficiency:} A $4.7\times$ lower end-to-end latency and $6.3\times$ lower memory footprint compared to the industry-standard Falco.
    \item \textbf{Formal Isolation:} A formally analyzed zero-copy channel that is resilient to TOCTOU attacks and a Wasm-based sandbox that blocks 100\% of rule-escape attempts.
\end{enumerate}

The remainder of this paper is organized as follows.
Section~\ref{sec:background} provides background on eBPF and Wasm.
Section~\ref{sec:threat-model} formalizes our threat model and security goals.
Section~\ref{sec:design} presents the system architecture and zero-copy mechanism.
Section~\ref{sec:implementation} details the implementation, followed by our extensive evaluation in Section~\ref{sec:evaluation}.
Section~\ref{sec:discussion} analyzes security and limitations, Section~\ref{sec:related-work} discusses related work, and Section~\ref{sec:conclusion} concludes.
